\begin{sepframe}{Functional code}
    {}
\end{sepframe}

\begin{frame}
    \frametitle{Functional code}
    \framesubtitle{Tips and tricks}

    \begin{itemize}[<+->]
        \item Make your objects \textbf{final},
        \item Make your properties \textbf{private} and \textbf{immutable},
        \item Make sure your functions are \textbf{total} and \textbf{pure}.
    \end{itemize}

    \note[item]{
        Some of you are getting to know me by now, and I couldn't make a
        presentation without discussing about functional code and how to
        improve your code with simple and basic tips like...
    }

    \note[item]{
        Since the last presentation we did and from the feedback that I had from
        colleagues, I noticed that the notion of final class is sometimes
        misunderstood. I'm going to show with an example why it is important to
        make sure that we can't extend classes.
    }

    \note[item]{
        Immutability is the trending word in I.T. since a couple of years now,
        and especially these last years. Its origin comes from
        the code and if we dig deeper, we will discover that it has roots in
        Lambda Calculus by Alonzo Church in 1930.
        Needless to say that making variables and objects immutable increase the
        ease of debugging and it is prone to the use of good design patterns.
    }

    \note[item]{
        Total functions are basically the gist of all of this including the
        reason why we have complexity and monads. I'm going to formalize the
        definition of them with a definition and some examples.
        I will also show what partial and pure functions are.
    }

\end{frame}

\begin{frame}
    \frametitle{Functional code}
    \framesubtitle{Tips and tricks}

\lstinputlisting[caption={Class is not final}]{src/session/functional/resources/immutability1-a.php}
\pause
\lstinputlisting[caption={When extending a class which is not final...}]{src/session/functional/resources/immutability1-b.php}
\end{frame}

\begin{frame}
    \frametitle{Functional code}
    \framesubtitle{Tips and tricks}

\lstinputlisting[caption={Class is not final}]{src/session/functional/resources/immutability2-a.php}
\pause
\lstinputlisting[caption={When extending a class which is not final...}]{src/session/functional/resources/immutability2-b.php}
\end{frame}

\begin{frame}
    \frametitle{Functional code}
    \framesubtitle{Tips and tricks}

\lstinputlisting[caption={Public properties}]{src/session/functional/resources/properties1-a.php}
\pause
\lstinputlisting[caption={Private properties}]{src/session/functional/resources/properties1-b.php}
\end{frame}

\begin{frame}
    \frametitle{Functional code}
    \framesubtitle{Tips and tricks}

    A total function is a function which is defined for all inputs.

    \pause

    \begin{definition}[Total function]
        \begin{align*}
        & f: A \mapsto B\\
        & \forall a \in A \implies \exists f(a) \in B
        \end{align*}
    \end{definition}
\end{frame}

\begin{frame}
    \frametitle{Functional code}
    \framesubtitle{Tips and tricks}

    A partial function is a function which is not defined for all inputs.

    \pause

    \begin{definition}[Partial function]
        \begin{align*}
        & f: A \rightharpoonup B\\
        & \exists x \in A \implies \not\exists f(x) \in B
        \end{align*}
    \end{definition}
\end{frame}

\begin{frame}
    \frametitle{Functional code}
    \framesubtitle{Tips and tricks}

    \begin{definition}[Total function]
        \begin{align*}
        & f: A \mapsto B\\
        & \forall a \in A \implies \exists f(a) \in B
        \end{align*}
    \end{definition}

    \begin{definition}[Partial function]
        \begin{align*}
        & f: A \mapsto B\\
        & \exists x \in A \implies \not\exists f(x) \in B
        \end{align*}
    \end{definition}
\end{frame}

\begin{frame}
    % TODO: Total functions are definitely less trouble makers than partials.
    \frametitle{Functional code}
    \framesubtitle{Tips and tricks}

    \begin{itemize}
        \item Total:\pause
        \\\lstinline[language=PHP]!function add (int $a, int $b): int!
        \pause
        \item Partial:\pause
        \\\lstinline[language=PHP]!function divide(int $a, int $b): float!
    \end{itemize}
\end{frame}

\begin{frame}
    \Huge
    \begin{center}
        What\pause\ the\pause\ link\pause\ ????
    \end{center}

    \note[item]{
        Some of you might say... but what is this guy saying? Why is he telling
        us all these theoratical stuff that I vaguely remember from school ?!
    }
    \note[item]{
        Don't worry, everything will come at once, everything is actually linked
        together.
    }
    \note[item]{
        Let's keep that in mind an move on.
    }
\end{frame}

\begin{frame}
    \frametitle{Functional code}
    \framesubtitle{Handling errors}

    \lstinputlisting[caption={A Doctrine repository definition}]{src/session/functional/resources/doctrine-repository1-a.php}
    \pause
    \lstinputlisting[caption={A Doctrine repository in use}]{src/session/functional/resources/doctrine-repository1-b.php}

    \note[item]{
        If you use Doctrine properly, this is the kind of stuff that you have
        all over your code.
        I don't know what you think, but this is usually where issues are
        coming. Why ?
        Simply because of the use of the null value. If the method find would
        not return null, then this check is definitely useless.
    }
    \note[item]{
        Let's use another example not involving Doctrine.
    }
\end{frame}

\begin{frame}
    \frametitle{Functional code}
    \framesubtitle{Tips and tricks}

    \lstinputlisting[caption={Return null}]{src/session/functional/resources/square-root1-a.php}
    \pause
    \lstinputlisting[caption={Throw exception}]{src/session/functional/resources/square-root1-b.php}
    \pause
    \lstinputlisting[caption={Return sentinel}]{src/session/functional/resources/square-root1-c.php}

    \note[item]{
        This idea here is to create a safeSquareRoot function which would return
        the square root of any numeric value, and we have to deal with negative
        numbers. And how should we deal the negative values? What are the
        options?
    }
    \note[item]{
        The first option we have is to return the amazing null value when the
        input value is a negative number.
    }
    \note[item]{
        The second option is to throw an exception and potentially stop the
        program.
    }
    \note[item]{
        The third option is to return a sentinel value.
    }
\end{frame}

\begin{frame}
    \frametitle{Functional code}
    \framesubtitle{Tips and tricks}

    \begin{itemize}
        \item Return null?\pause
            \\\textcolor{ecgrey!50}{Mixing types implies more code for the end
            user, prevent composition.}
        \item Throw exceptions?\pause
            \\\textcolor{ecgrey!50}{Great power means great responsibility, has
            the power to stop the program.}
        \item Return sentinel?\pause
            \\\textcolor{ecgrey!50}{Prone to incertitude when it returns 0.}
    \end{itemize}

    \note[item]{
        Return null ... explain and give an example with Doctrine.
    }
    \note[item]{
        Exceptions are great, but needs to be correctly documented and handled.
    }
    \note[item]{
        Sentinels, avoid that, it's an old thing and extremelly prone to
        incertitude.
    }
\end{frame}

\begin{frame}
    \frametitle{Functional code}
    \framesubtitle{Return null?}

    \begin{itemize}[<+->]
        \item Universal error value that doesn't carry any relevant information
            in case of issues.
        \item Widespread, implemented in most languages
        \item Considered as a code smell and a "the billion-dollar mistake"
        \item Prevent composition pattern
    \end{itemize}

    \note[item]{
        Java, C,... implements nullable objects BY DEFAULT, controversial.
    }
\end{frame}

\begin{frame}
    \frametitle{Functional code}
    \framesubtitle{Throw exceptions?}

    \begin{itemize}[<+->]
        \item Indicate something unforeseen by the developer
        \item Expensive, destructive
        \item Definitely prevent composition pattern
    \end{itemize}

    \note[item]{
        Should be carefully used.
    }
\end{frame}

\begin{frame}
    \frametitle{Functional code}
    \framesubtitle{Return sentinel?}

    \begin{itemize}[<+->]
        \item Potentially prevent composition pattern
        \item Must be carefully done, each type has its own unit type
        \begin{itemize}[<+->]
            \item \lstinline[language=PHP]!function divide(int $a, int $b): float|string!
            \item \lstinline[language=PHP]!function divide(int $a, int $b): float|array!
            \item \lstinline[language=PHP]!function divide(int $a, int $b): float|null!
            \item \lstinline[language=PHP]!function divide(int $a, int $b): float|int!
        \end{itemize}
        \item In PHP, this is sadly happening
    \end{itemize}

    \note[item]{
        Ask the audience in which PHP function this happens ? (strpos)
    }
\end{frame}

\begin{frame}
    \frametitle{Functional code}
    \framesubtitle{What is our best option?}

    \begin{itemize}[<+->]
        \item Not using PHP ?
        \item Avoid using strict types ?
        \item Wrap the incertitude in a well-know object?
    \end{itemize}

    \note[item]{
        Take a break, let's make a summary of what we've seen (TODO)
    }
    \note[item]{
        Start to slowly introduce the concept of monads using Schrodinger's cat.
    }
\end{frame}

\begin{frame}
    \begin{figure}
        \begin{overprint}
        \onslide<1>
            \begin{center}
                \includegraphics[width=8em]{src/session/functional/resources/cat.png}
            \end{center}
        \onslide<2->
            \begin{center}
                \includegraphics[width=10em]{src/session/functional/resources/box.png}
            \end{center}
        \end{overprint}
    \end{figure}

    \pause
    \pause

    How about the idea of abstracting the incertitude by creating an object that
    would wrap a value in a well-know object?\\

    \pause

    \vspace{1em}
    An object that can be \textbf{maybe} null or \textbf{maybe} the expected
    value?
\end{frame}

\begin{frame}
    \lstinputlisting[caption={Draft Maybe class}]{src/session/functional/resources/maybe-draft1-a.php}
\end{frame}

\begin{frame}
    \lstinputlisting[caption={Draft Maybe class}]{src/session/functional/resources/maybe-draft1-b.php}
\end{frame}

\begin{frame}
    \lstinputlisting[caption={Draft Maybe class}]{src/session/functional/resources/maybe-draft1-c.php}
\end{frame}

\begin{frame}
    \lstinputlisting[caption={Draft Maybe class}]{src/session/functional/resources/maybe-draft1-d.php}
    % Problem is still there, we have to write pure function and
    % avoid mixed types.
\end{frame}

\begin{frame}
    \lstinputlisting[caption={Draft Maybe class}]{src/session/functional/resources/maybe-draft1-e.php}
    % Problem is still there, we have to write pure function and
    % avoid mixed types.
\end{frame}

\begin{frame}
    \lstinputlisting[caption={Draft Maybe class}]{src/session/functional/resources/maybe-draft1-f.php}
    % Problem is still there, we have to write pure function and
    % avoid mixed types.
\end{frame}

\begin{frame}
    \frametitle{Functional code}
    \framesubtitle{The Maybe monad}

    \begin{itemize}[<+->]
        \item Focus on the happy scenario
        \item Return null when something goes wrong
        \item Unable to customize the value when something is wrong
        \item Back to square 1.
    \end{itemize}
\end{frame}

\begin{frame}
    Maybe we could customize the return value in case of something wrong happen?\\
    \vspace{1em}
    \pause
    While still being able to decide what to do when an error happen?
\end{frame}

\begin{frame}
    \lstinputlisting[caption={Draft Either class}]{src/session/functional/resources/either-draft1-f.php}
\end{frame}

\begin{frame}
    \centering
    \Huge So these are monads !
\end{frame}

\begin{frame}[plain]
    \centering
    \vskip 0pt plus 1filll
    \Huge Should I use monads ?
    \vskip 0pt plus 1filll
    \makebox[\linewidth]{\includegraphics[width=\paperwidth]{src/session/functional/resources/iceberg-top.png}}
\end{frame}

\begin{frame}[plain]
    \makebox[\linewidth]{\includegraphics[width=\paperwidth]{src/session/functional/resources/iceberg-bottom.png}}
    \centering
    \vskip 0pt plus 1filll
    \Large Technically, a monad is a monoid in the monoidal category of
    endofunctors equipped with functor composition as its product.
    \vskip 0pt plus 1filll
\end{frame}

% Silver lining
\begin{frame}[plain]
    \makebox[\linewidth]{\includegraphics[width=\paperwidth]{src/session/functional/resources/iceberg.png}}
\end{frame}

\begin{frame}
    \begin{itemize}
        \item Maybe: Deal with a value and null
        \item Either: Deal with a value and another value
        \item ...Reader, Writer, IO, List,... There are many other monads these
            two are just the tip of the iceberg.
        \item Monads may have different names in languages (Optional, Result, Future, etc etc)
        \item Monads are the design pattern of functional languages
        \item Underlying theory can be very complex
        \item 1988, Eugenio Moggi
        \item We don't need to understand them to use them !
    \end{itemize}
\end{frame}

\begin{frame}
    \centering
    \Huge Should I use monads ?
\end{frame}

\begin{frame}
    \begin{itemize}
        \item Monads can abstract most of your function binding, turning a lot of
            noise into magic.
        \item It modularizes logical unit of computation and facilitate refactoring
        \item Reduce complexity of the code, abstracting the logic into one well
            known object that does only thing.
        \item Either: Deal with a value and another value
        \item For educational purposes ? Demystifing monads can help understand
            the benefits and perks of functional programming, and make you a better
            developer.
    \end{itemize}
\end{frame}
